\chapter{General Introduction}
\label{chap:introduction}

The MFJA robots, transport cells and cameras did not yet form a coherent
simulation for reliable planning and experiments. The original assignment was
to build the Room~315 environment and reusable launch foundation. Its early
completion enabled Phase~2 on task-level English and paired-camera observations.
PlanSys2 requests symbolic plans from rules written in the Planning Domain
Definition Language (PDDL) \citep{plansys2}. PDDL defines the objects, state,
actions, preconditions and effects \citep{pddl}. Learned models supply checked
task and scene descriptions; they never issue commands.

\paragraph{Problem and objectives.}

I used Autodesk Inventor for the computer-aided-design (CAD) geometry
\citep{autodeskinventor}, Gazebo to execute the scenes
\citep{gazeboharmonic}, and ROS~2 to connect their software components
\citep{ros2jazzyinterfaces}.

\noindent\textbf{Engineering problem.} \emph{How can the heterogeneous robots,
rail cells and overhead-camera observations of Room~315 be integrated into a
reproducible Gazebo environment that transforms constrained English requests
into checked shuttle motion without giving learned models direct command
authority?}

To answer this question, the project pursued four engineering objectives:

\begin{enumerate}[label=\textbf{O\arabic*.},leftmargin=*,itemsep=0.22em]
  \item design the third-floor geometry and furnished-room contents, then
  integrate the Room~315 robots and cells in a maintainable Gazebo environment
  with selectable launch and unified ROS~2 access;
  \item implement calibrated, directed multi-shuttle transport with device
  semantics, routing, reservations and separation constraints;
  \item define typed command and state interfaces, together with supervision
  rules for checking commands and confirming their effects; and
  \item design and evaluate a neuro-symbolic workflow that combines
  natural-language goal interpretation and paired-camera state estimation with
  PDDL/PlanSys2 planning and software-supervised execution.
\end{enumerate}

O1--O4 progress from simulation to task-level interaction and are assessed in
\cref{tab:objective-results}.

\paragraph{Method.}

Development followed short build--test--review cycles. I checked each layer
before integrating it, then froze the software, data and model versions used for
evaluation (\cref{app:reproducibility}). Preserving consistent geometry,
identity and state across layers was the main integration challenge.

The delivered Gazebo system starts six isolated robot instances and coordinates
eight shuttles. The frozen estimator reached 99.94\% segment accuracy; 12/12
closed-loop tasks and 5/5 fault cases met their predefined outcomes
(\cref{chap:results}).

\begin{keybox}[Project demonstration videos]
Project demonstrations are available in the
\href{https://1drv.ms/f/c/b9a744d4252a8ea4/IgDZe0nBqhLRQYNN5V2g-S-nAVb9ZIn0VUSxfA4CWJeYxTA?e=eGmKr2}
{\textcolor{mfjablue}{\textbf{\underline{online demonstration-video folder}}}}.
The videos provide a visual
overview; the measured results and reproducibility records remain in
\cref{chap:results,app:reproducibility}.
\end{keybox}

\paragraph{Report structure.}
Chapters~2--4 introduce the host, mission and design choices. Chapters~5--8
develop the simulation, robot, transport and interface foundations; Chapters~9
and~10 add the AI inputs and checked planning loop. Chapters~11--13 assess the
evidence, limitations and contributions against O1--O4.
